\chapter{电子碰撞、束流统计与费米关联}\label{chap:electron-collision-dynamics}

电子可以交换一个玻色激发，也可以直接改变一个束缚靶的内部态。
若出射电子不被记录，靶的演化就由联合幺正算符的偏迹决定。
把这个单次作用重复起来，还要规定电子是否独立、何时到达，以及电子波包重叠时怎样保持费米反对称性。
这些条件决定单电子相干能否驱动靶、随机束流是否可视为无记忆储库，以及两电子关联应怎样解释。

\section{从联合碰撞到靶信道}\label{sec:collision-channel-beam-dp19}\label{subsec:single-electron-collision-dp19}

束缚靶的离散内部态可以直接与电子连续态耦合，无需先改写成玻色谐振子。保留电子的初末动量核，才能在计算后判断转移动量、靶态截断及独立碰撞的条件。

取任意有限维靶 Hilbert 空间
\begin{equation*}
\mathcal H_A=\operatorname{span}\{|i\rangle\}_{i=1}^{N},
\qquad H_A=\sum_iE_i|i\rangle\langle i|,
\qquad \omega_{ij}=\frac{E_i-E_j}{\hbar}.
\end{equation*}
在已固定的正能自旋通道内，用 $\kappa=(p-p_0)/\hbar$ 标记自由电子纵向动量偏移。在电子数守恒、且相互作用对电子场只含一产生一湮灭算符的一体耦合范围内，初、末两个连续动量都应先保留：
\begin{equation*}
H_{eA}
=\sum_{ij}\int\dd \kappa\,\dd \kappa'\,
V_{ij}(\kappa',\kappa)c_{\kappa'}^\dagger c_\kappa\,|i\rangle\langle j|,
\qquad
V_{ji}(\kappa,\kappa')=V_{ij}^*(\kappa',\kappa).
\end{equation*}
在高能准直、电子带宽窄且矩阵元主要依赖转移动量 $u=\kappa-\kappa'$ 时，进一步定义电子平移算符
\begin{equation*}
b_u=\int\dd \kappa\,c_\kappa^\dagger c_{\kappa+u}.
\end{equation*}
设实际电子支持满足 $|\kappa|,|\kappa'|\le\Delta\kappa_e$，并选取只依赖 $u$ 的核 $V_{ij}(u)$。以代表性相互作用核尺度 $V_*$ 定义
\begin{equation*}
 \epsilon_V
 =\frac{1}{V_*}
 \sup_{i,j;\kappa,\kappa'}
 |V_{ij}(\kappa',\kappa)-V_{ij}(\kappa-\kappa')|\ll1.
\end{equation*}
再令 $T_c=t_f-t_i$ 为碰撞作用时间，并以
\begin{equation*}
 M_2=\sup_{|p-p_0|\le\hbar\Delta\kappa_e}|E''(p)|,
 \qquad
 \epsilon_{\rm disp}
 =\frac{T_c}{2\hbar}M_2(\hbar\Delta\kappa_e)^2\ll1
\end{equation*}
控制自由色散在线性项之后累积的最大二阶相位。于是
\begin{equation*}
 H_{eA}
 =\sum_{ij}\int\dd u\,V_{ij}(u)b_u|i\rangle\langle j|
 +R_V,
\end{equation*}
而顶点相位在 $O(\epsilon_{\rm disp})$ 精度内为 $\ee^{\ii(\omega_{ij}-v_0u)t}$。有限碰撞时间由 $|\omega_{ij}-v_0u|T_c\lesssim1$ 选择可分辨转移动量；$T_c\to\infty$ 时才收敛到 $u=\omega_{ij}/v_0$ 的尖锐共振。

核的点态相对误差还必须乘上工作窗大小，才能成为传播误差。
\begin{proposition}[有限动量窗内的转移核近似]
在 $L^2([-\Delta\kappa_e,\Delta\kappa_e])\otimes\mathbb C^N$ 上，
若各块余项核满足 $|R_{V,ij}(\kappa',\kappa)|\le V_*\epsilon_V$，则
\[
\|R_V\|\le2N\Delta\kappa_e V_*\epsilon_V.
\]
将两个模型限制在同一工作空间，并只改变上述核及自由色散的二阶余项时，
\begin{equation}
\|U_c-U_{\rm tr}\|
\le\frac{2N T_c\Delta\kappa_e V_*}{\hbar}\epsilon_V+\epsilon_{\rm disp}.
\label{eq:collision-transfer-kernel-error}
\end{equation}
$V_*$ 具有能量乘长度的单位；在完整连续态空间使用该界，还须独立计入工作窗外的振幅。
\end{proposition}
\begin{proof}
核作为积分算符作用时，每行与每列的绝对积分和均不超过
$2N\Delta\kappa_e V_*\epsilon_V$。更具体地，对各输入分量应用 Cauchy--Schwarz，
再对 $i$ 求和、对输出动量积分，可得
\[
\|R_Vf\|^2\le(2N\Delta\kappa_e V_*\epsilon_V)^2\|f\|^2.
\]
色散 Taylor 余项在窗内的乘法算符范数不超过
$M_2(\hbar\Delta\kappa_e)^2/2$。
将两个 Hamiltonian 的差代入 Duhamel 积分，并利用两传播子的幺正性，
便得到相互作用时长乘这个总能量误差除以 $\hbar$ 的上界。
所以 $\epsilon_V\ll1$ 是核形状的比较，实际动力学还要检验\eqnrefs{eq:collision-transfer-kernel-error}。
\end{proof}

一次电子通过靶系统的精确作用由时间序幺正算符
\begin{equation}
\boxed{
U_c=\mathcal T\exp\!\left[-\frac{\ii}{\hbar}\int_{t_i}^{t_f}\dd t\,V_I(t)\right].}
\label{eq:general_collision_exact_unitary}
\end{equation}
给出。若入射电子密度矩阵为 $\rho_e$，离开后电子不再被追踪，则靶系统经历精确 CPTP 碰撞信道
\begin{equation}
\boxed{
\mathscr E_e(\rho_A)
=\operatorname{Tr}_e[U_c(\rho_e\otimes\rho_A)U_c^\dagger].}
\label{eq:general_electron_target_cptp_map}
\end{equation}
若把电子初态谱分解为 $\rho_e=\sum_\mu p_\mu|\mu\rangle\langle\mu|$，并在电子输出空间选完备基 $|r\rangle$，则同一信道等价地写成
\begin{equation}
\boxed{
K_{r\mu}=\sqrt{p_\mu}\,\langle r|U_c|\mu\rangle,
\qquad
\mathscr E_e(\rho_A)=\sum_{r\mu}K_{r\mu}\rho_AK_{r\mu}^\dagger,
\qquad
\sum_{r\mu}K_{r\mu}^\dagger K_{r\mu}=\mathbb I_A.}
\label{eq:general_collision_kraus_representation}
\end{equation}
最后一个完备关系说明单次碰撞信道严格保迹；Kraus 和式同时保证完全正性。

\begin{theorem}[任意电子初态诱导的靶 Kraus 信道]
设电子初态 $\rho_e=\sum_\mu p_\mu|\mu\rangle\langle\mu|$ 与任意靶输入态因子化，联合碰撞 $U_c$ 幺正。靶偏迹映射为 CPTP，Kraus 算符 $K_{r\mu}=\sqrt{p_\mu}\langle r|U_c|\mu\rangle_e$ 满足\eqnrefs{eq:general_collision_kraus_representation}。电子连续谱时以完备测度积分取代末态求和。
\end{theorem}
\begin{proof}
先明确这里每个算符作用在哪个空间。$U_c$ 作用在联合空间 $\mathcal H_e\otimes\mathcal H_A$；固定电子初、末态以后，偏矩阵元
\[
{}_e\!\langle r|U_c|\mu\rangle_e
\]
仍然是作用在靶空间 $\mathcal H_A$ 上的算符，而不是一个复数。正是这些偏矩阵元构成靶系统的 Kraus 算符。

把任意电子初态作谱分解
\[
\rho_e=\sum_\mu p_\mu|\mu\rangle\langle\mu|,
\qquad p_\mu\ge0,
\qquad \sum_\mu p_\mu=1,
\]
并在电子输出空间插入完备关系
\[
\sum_r|r\rangle\langle r|=\mathbb I_e.
\]
由偏迹的定义，
\begin{align*}
\mathscr E_e(\rho_A)
&=\sum_r{}_e\!\langle r|
U_c(\rho_e\otimes\rho_A)U_c^\dagger
|r\rangle_e\\
&=\sum_{r\mu}p_\mu
{}_e\!\langle r|U_c|\mu\rangle_e\,
\rho_A\,
{}_e\!\langle\mu|U_c^\dagger|r\rangle_e.
\end{align*}
定义
\[
K_{r\mu}\equiv\sqrt{p_\mu}
{}_e\!\langle r|U_c|\mu\rangle_e,
\]
便得到算符和式
\[
\mathscr E_e(\rho_A)
=\sum_{r\mu}K_{r\mu}\rho_AK_{r\mu}^\dagger.
\]

接着验证保迹性，而不是只引用定理。对靶空间上的任意向量 $|\varphi\rangle$，
\begin{align*}
\left\langle\varphi\left|
\sum_{r\mu}K_{r\mu}^\dagger K_{r\mu}
\right|\varphi\right\rangle
&=\sum_\mu p_\mu\sum_r
\left\|
{}_e\!\langle r|U_c(|\mu\rangle_e\otimes|\varphi\rangle_A)
\right\|_A^2\\
&=\sum_\mu p_\mu
\left\|U_c(|\mu\rangle_e\otimes|\varphi\rangle_A)\right\|_{eA}^2\\
&=\sum_\mu p_\mu\|\varphi\|_A^2
=\|\varphi\|_A^2.
\end{align*}
第二行使用电子输出基的 Parseval 完备性，第三行使用 $U_c$ 的幺正性。因为该二次型对任意 $|\varphi\rangle$ 都成立，故
\[
\sum_{r\mu}K_{r\mu}^\dagger K_{r\mu}=\mathbb I_A.
\]
于是
\begin{align*}
\Tr_A\mathscr E_e(\rho_A)
&=\sum_{r\mu}\Tr_A(K_{r\mu}\rho_AK_{r\mu}^\dagger)\\
&=\Tr_A\!\left[
\rho_A\sum_{r\mu}K_{r\mu}^\dagger K_{r\mu}
\right]
=\Tr_A\rho_A.
\end{align*}

最后验证“完全正”中“完全”的含义。取一个任意维度、完全不参与碰撞的参考系统 $R$，对任意正算符 $X_{RA}\ge0$，有
\begin{equation*}
(\mathbb I_R\otimes\mathscr E_e)(X_{RA})
=\sum_{r\mu}
(\mathbb I_R\otimes K_{r\mu})
X_{RA}
(\mathbb I_R\otimes K_{r\mu}^\dagger).
\end{equation*}
对任意 $|\Xi\rangle\in\mathcal H_R\otimes\mathcal H_A$，
\begin{align*}
\langle\Xi|(\mathbb I_R\otimes\mathscr E_e)(X_{RA})|\Xi\rangle
=\sum_{r\mu}
\langle\Xi_{r\mu}|X_{RA}|\Xi_{r\mu}\rangle\ge0,
\end{align*}
其中 $|\Xi_{r\mu}\rangle=(\mathbb I_R\otimes K_{r\mu}^\dagger)|\Xi\rangle$。因此即使靶系统与任意外部参考系统纠缠，碰撞后联合算符仍保持正性；这就是完全正性。

于是正文\eqnrefs{eq:general_collision_kraus_representation} 的三个式子并非三条独立假设：Kraus 算符来自联合幺正算符的电子偏矩阵元，完备关系来自幺正性，而完全正性来自同一 Kraus 和式在任意参考扩张上的正性。
\end{proof}


弱碰撞的小量必须由实际时间积分定义：
\begin{equation*}
\epsilon_{\rm coll}
=\frac1\hbar\int_{t_i}^{t_f}\dd t\,\|V_I(t)\|\ll1
\end{equation*}
（无界相互作用时应改用能量约束或投影后的算符范数）。Magnus 展开到二阶写成
\begin{equation}
\boxed{
U_c
=\exp[-\ii(V_1+V_2)+O(\epsilon_{\rm coll}^3)],
\qquad
V_1=\frac1\hbar\int\dd t\,V_I(t),}
\label{eq:general_collision_dyson_magnus_second_order}
\end{equation}
其中二阶时间序生成元
\begin{equation}
\boxed{
V_2
=-\frac{\ii}{2\hbar^2}
\int_{t_i}^{t_f}\dd t_1\int_{t_i}^{t_1}\dd t_2\,
[V_I(t_1),V_I(t_2)].}
\label{eq:general_collision_second_magnus_generator}
\end{equation}
若使用瞬时冲量模型，还须由具体包络证明这项为零、已计入有效参数，或小于要求的误差。

\begin{proposition}[有界碰撞传播子的二阶 Magnus 生成元]
设 $V_I(t)$ 自伴且有界，$\epsilon_{\rm coll}=\hbar^{-1}\int\|V_I(t)\|\dd t\le1/4$。
传播子的解析对数前两项为\eqnrefs{eq:general_collision_dyson_magnus_second_order}与\eqnrefs{eq:general_collision_second_magnus_generator}，
余项满足 $\|R_\Omega^{(3)}\|\le2\epsilon_{\rm coll}^3$。时间交换子未消失时，二阶项不能删去。
\end{proposition}
\begin{proof}
Magnus 展开不是把 Dyson 级数“重新命名”为指数，而是要求一个指数的逐阶展开与时间排序级数逐项相等。先把精确传播子写到二阶：
\begin{align*}
 U_c
 &=\mathbb I+D_1+D_2+R_D^{(3)},\\
 D_1
 &=-\frac{\ii}{\hbar}\int_{t_i}^{t_f}\dd t_1\,V_I(t_1),\\
 D_2
 &=-\frac{1}{\hbar^2}
 \int_{t_i}^{t_f}\dd t_1
 \int_{t_i}^{t_1}\dd t_2\,
 V_I(t_1)V_I(t_2).
\end{align*}
现在假设同一个传播子可写为
\[
 U_c=\exp(\Omega_1+\Omega_2+R_\Omega^{(3)}),
\]
其中 $\Omega_n$ 是相互作用的 $n$ 阶量。把指数只展开到二阶，
\begin{equation*}
 U_c
 =\mathbb I+\Omega_1+\Omega_2
 +\frac12\Omega_1^2+R^{(3)}.
\end{equation*}
与 Dyson 展开逐阶比较立刻得到
\begin{equation*}
 \Omega_1=D_1,
 \qquad
 \Omega_2=D_2-\frac12D_1^2.
\end{equation*}
关键数学步骤在第二式。平方的一阶积分覆盖整个正方形时间域：
\begin{align*}
 \frac12D_1^2
 &=-\frac{1}{2\hbar^2}
 \int_{t_i}^{t_f}\dd t_1\int_{t_i}^{t_f}\dd t_2\,
 V_I(t_1)V_I(t_2)\\
 &=-\frac{1}{2\hbar^2}
 \int_{t_1>t_2}\dd t_1\dd t_2
 \left[V_I(t_1)V_I(t_2)+V_I(t_2)V_I(t_1)\right],
\end{align*}
其中第二行把正方形分成 $t_1>t_2$ 与 $t_2>t_1$ 两个三角形，并在后一块交换积分变量。另一方面
\begin{equation*}
 D_2=-\frac{1}{\hbar^2}
 \int_{t_1>t_2}\dd t_1\dd t_2\,
 V_I(t_1)V_I(t_2).
\end{equation*}
相减后，对称乘积部分消失，只剩交换子：
\begin{align*}
 \Omega_2
 &=-\frac{1}{2\hbar^2}
 \int_{t_1>t_2}\dd t_1\dd t_2\,
 \left[V_I(t_1),V_I(t_2)\right].
\end{align*}
因此若定义无量纲厄米生成元
\begin{align*}
 V_1&\equiv\frac1\hbar\int_{t_i}^{t_f}\dd t\,V_I(t),\\
 V_2&\equiv-\frac{\ii}{2\hbar^2}
 \int_{t_1>t_2}\dd t_1\dd t_2\,
 [V_I(t_1),V_I(t_2)],
\end{align*}
便有 $\Omega_1=-\ii V_1$、$\Omega_2=-\ii V_2$，从而
\begin{equation*}
 U_c
 =\exp\!\left[-\ii(V_1+V_2)+R_\Omega^{(3)}\right].
\end{equation*}
还需控制余项。记 $\epsilon=\epsilon_{\rm coll}$，$X=U_c-I$，
$x=\ee^\epsilon-1$、$y=\ee^\epsilon-1-\epsilon$。
Dyson 级数给 $\|X\|\le x<1$，$\|X-D_1\|\le y$，
$\|R_D^{(3)}\|\le\ee^\epsilon\epsilon^3/6$。
因此 $\log(I+X)=X-X^2/2+R_{\log}$ 收敛，且
\[
\|R_{\log}\|\le\frac{x^3}{3(1-x)},\qquad
\|X^2-D_1^2\|\le(x+\epsilon)y.
\]
把这些界与 $\Omega_1+\Omega_2=D_1+D_2-D_1^2/2$ 比较，得
\[
\|R_\Omega^{(3)}\|\le\frac{\ee^\epsilon\epsilon^3}{6}
+\frac{(x+\epsilon)y}{2}+\frac{x^3}{3(1-x)}\le2\epsilon^3.
\]
最后一步在 $\epsilon\le1/4$ 使用
$x\le\ee^{1/4}\epsilon$、$y\le\ee^{1/4}\epsilon^2/2$、$1-x\ge2-\ee^{1/4}$。
所以脉冲只由累计面积 $V_1$ 决定，需要不同时刻的作用对易；否则时间先后顺序本身贡献 $V_2$。
\end{proof}


\section{束流的一体相干与费米涨落}

整形电子对靶系统的作用通过电子平移相干性进入。对仅含转移的模型定义
\begin{equation*}
I_e(u)=\operatorname{Tr}_e[\rho_e b_u]
=\int\dd \kappa\,\psi_e(\kappa+u)\psi_e^*(\kappa)
\end{equation*}
（第二式只用于单电子纯态；多电子束流应使用其一体密度矩阵）。它是电子一体密度矩阵在固定动量差上的切片：$|I_e|$ 控制两条靶跃迁路径能够保留多少相干，$\arg I_e$ 提供相对驱动相位。电子的相位调制、自由色散或其他纵向整形制备可调的 $I_e(u)$；碰撞信道则以该入射相干函数作为输入。

单电子碰撞信道还隐含一个经常被忽略的统计前提：不同入射电子可以被视为彼此独立的辅助系统。对相干调制电子列、简并电子气或具有显著二电子关联的束流，这个假设一般不成立。此时靶跃迁首先读取的不是“平均电子数”，而是与目标频率匹配的电子一体算符的\emph{两时关联函数}。

设某个目标跃迁 $|a\rangle\to|b\rangle$ 的 Bohr 频率为 $\omega_{ba}>0$，并把与该跃迁矩阵元收缩后的电子一体驱动算符记为
\begin{equation}
\boxed{
\hat F(t)
=\sum_{\alpha\beta}F_{\alpha\beta}(t)
\hat c_\alpha^\dagger\hat c_\beta .}
\label{eq:fermionic_beam_onebody_drive_operator}
\end{equation}
这里 $\alpha,\beta$ 可以包含连续动量、自旋和横向模；求和在连续指标上理解为相应积分。弱跃迁极限中，对末态电子全部求和后，靶的二阶激发概率为
\begin{equation}
\boxed{
P_{a\to b}^{(2)}
=\frac{1}{\hbar^2}
\int_{t_i}^{t_f}\dd t\int_{t_i}^{t_f}\dd t'\,
\ee^{+\ii\omega_{ba}(t-t')}
\left\langle
\hat F^\dagger(t')\hat F(t)
\right\rangle_e .}
\label{eq:target_excitation_fermionic_beam_correlator}
\end{equation}
因此目标系统实际上充当了一个频率选择的电子关联函数探针。

\begin{proposition}[束流驱动的二阶靶激发概率]
设初态为 $|a\rangle\langle a|\otimes\rho_e$，
相关跃迁项是 $|b\rangle\langle a|F(t)\ee^{\ii\omega_{ba}t}$ 及其厄米共轭，
$F(t)$ 的单位为能量。一级振幅给出的概率为\eqnrefs{eq:target_excitation_fermionic_beam_correlator}。
令 $B=\int_{t_i}^{t_f}\ee^{\ii\omega_{ba}t}F(t)\dd t$，则该概率也等于
$\hbar^{-2}\|B\rho_e^{1/2}\|_{\rm HS}^2\ge0$。
若完整相互作用有界且 $\epsilon=\int\|V_I\|\dd t/\hbar\ll1$，
一般多能级问题中所舍弃的概率项为 $O(\epsilon^3)$；无界情形需要初态上的独立传播控制。
\end{proposition}
\begin{proof}
取所考察跃迁的相互作用项。取相互作用绘景中与 $a\to b$ 共振的部分
\[
\hat V_I(t)
=|b\rangle\langle a|\,
\hat F(t)\ee^{+\ii\omega_{ba}t}
+\mathrm{H.c.}
\]
并从 $|a\rangle\langle a|\otimes\rho_e$ 出发。一阶振幅对电子末态求模平方并完备求和，等价于密度矩阵二阶表达
\begin{align*}
P_{a\to b}^{(2)}
&=\frac1{\hbar^2}
\int\dd t\,\dd t'\,
\ee^{+\ii\omega_{ba}(t-t')}
\Tr_e[
\hat F(t)\rho_e\hat F^\dagger(t')]\\
&=\frac1{\hbar^2}
\int\dd t\,\dd t'\,
\ee^{+\ii\omega_{ba}(t-t')}
\langle\hat F^\dagger(t')\hat F(t)\rangle_e,
\end{align*}
即正文\eqnrefs{eq:target_excitation_fermionic_beam_correlator}。把刚刚证明的 Wick 分解代入，就同时得到相干 $N_e^2$ 项、量子涨落 $N_e$ 项与费米交换阻塞，而无需另行假定经典聚束因子。
两个时间积分合成 $B$ 后，循环使用迹便得
$\operatorname{Tr}(B\rho_eB^\dagger)/\hbar^2$
和所述 Hilbert--Schmidt 范数。
有界 Dyson 级数的一级振幅为 $O(\epsilon)$，余振幅为 $O(\epsilon^2)$；
它们的交叉项至多为 $O(\epsilon^3)$，给出概率截断阶数。
\end{proof}


对任意保持粒子数的 高斯费米态，定义一体密度矩阵
\begin{equation*}
\Gamma_{\alpha\beta}
\equiv
\langle \hat c_\beta^\dagger\hat c_\alpha\rangle,
\qquad
0\le\Gamma\le I .
\end{equation*}
Wick 定理把上式中的四费米平均严格化成
\begin{equation}
\boxed{
\begin{aligned}
\langle\hat F^\dagger(t')\hat F(t)\rangle
={}&
\Tr[F^\dagger(t')\Gamma]\,
\Tr[F(t)\Gamma]\\
&+\Tr\!\left[
F^\dagger(t')(I-\Gamma)F(t)\Gamma
\right].
\end{aligned}}
\label{eq:fermionic_beam_coherent_incoherent_split}
\end{equation}
第一项是平均一体驱动的相干平方；若许多电子在目标 Fourier 分量上锁相，它可以呈现 $N_e^2$ 型相干增强。第二项是不可由平均电流替代的量子涨落项，其中 $I-\Gamma$ 明确给出 Pauli 阻塞：一个已被占据的末电子模不能再次接受同一个费米子跃迁。对稀薄非简并束流，$\Gamma$ 的本征占据远小于一，第二项退化为 $\Tr(F^\dagger F\Gamma)$ 的独立电子贡献；对 Slater 行列式，$\Gamma^2=\Gamma$，交换反关联则被完整保留。

\begin{lemma}[准自由费米态的一体乘积恒等式]
设数守恒准自由费米态的一体密度为 $0\le\Gamma\le I$，$\hat X=c^\dagger Xc$、$\hat Y=c^\dagger Yc$，所有迹存在。则 $\langle\hat X\hat Y\rangle=\Tr(X\Gamma)\Tr(Y\Gamma)+\Tr[X(I-\Gamma)Y\Gamma]$。取 $X=F^\dagger,Y=F$ 得\eqnrefs{eq:fermionic_beam_coherent_incoherent_split}；一般相关费米态还包含四点连通项。
\end{lemma}
\begin{proof}
先从一个保持粒子数的一体算符
\[
\hat X=\hat c^\dagger X\hat c
=\sum_{ij}X_{ij}\hat c_i^\dagger\hat c_j
\]
开始。定义
\[
\Gamma_{ij}=\langle\hat c_j^\dagger\hat c_i\rangle,
\]
则
\[
\langle\hat X\rangle
=\sum_{ij}X_{ij}\Gamma_{ji}
=\Tr(X\Gamma).
\]
目标是计算两个一般一体算符
$\hat X=\hat c^\dagger X\hat c$ 与
$\hat Y=\hat c^\dagger Y\hat c$ 的乘积平均：
\[
\langle\hat X\hat Y\rangle
=\sum_{ijkl}X_{ij}Y_{kl}
\langle
\hat c_i^\dagger\hat c_j
\hat c_k^\dagger\hat c_l
\rangle.
\]

对保持粒子数的 高斯费米态，异常配对
$\langle cc\rangle$ 与 $\langle c^\dagger c^\dagger\rangle$ 为零。Wick 定理只留下两个非零配对：
\begin{align*}
\langle
\hat c_i^\dagger\hat c_j
\hat c_k^\dagger\hat c_l
\rangle
={}&
\langle\hat c_i^\dagger\hat c_j\rangle
\langle\hat c_k^\dagger\hat c_l\rangle\\
&+
\langle\hat c_i^\dagger\hat c_l\rangle
\langle\hat c_j\hat c_k^\dagger\rangle.
\end{align*}
反对易关系给
\[
\langle\hat c_j\hat c_k^\dagger\rangle
=\delta_{jk}-\langle\hat c_k^\dagger\hat c_j\rangle
=\delta_{jk}-\Gamma_{jk},
\]
所以
\[
\langle
\hat c_i^\dagger\hat c_j
\hat c_k^\dagger\hat c_l
\rangle
=
\Gamma_{ji}\Gamma_{lk}
+\Gamma_{li}(\delta_{jk}-\Gamma_{jk}).
\]

逐项代回。第一项为
\begin{align*}
\sum_{ijkl}X_{ij}Y_{kl}
\Gamma_{ji}\Gamma_{lk}
&=
\Tr(X\Gamma)\Tr(Y\Gamma).
\end{align*}
第二项中的 $\delta_{jk}$ 给
\[
\sum_{ijl}X_{ij}Y_{jl}\Gamma_{li}
=\Tr(XY\Gamma),
\]
而交换减项为
\[
\sum_{ijkl}X_{ij}Y_{kl}\Gamma_{li}\Gamma_{jk}
=\Tr(X\Gamma Y\Gamma).
\]
因此得到一般恒等式
\begin{equation*}
\boxed{
\langle\hat X\hat Y\rangle
=
\Tr(X\Gamma)\Tr(Y\Gamma)
+\Tr(XY\Gamma)
-\Tr(X\Gamma Y\Gamma).}
\end{equation*}
把最后两项合并：
\[
\Tr(XY\Gamma)-\Tr(X\Gamma Y\Gamma)
=
\Tr[X(I-\Gamma)Y\Gamma].
\]
取
\[
X=F^\dagger(t'),
\qquad
Y=F(t),
\]
则
\begin{align*}
\langle\hat F^\dagger(t')\hat F(t)\rangle
={}&\Tr[F^\dagger(t')\Gamma]\Tr[F(t)\Gamma]\\
&+\Tr\!\left[F^\dagger(t')(I-\Gamma)F(t)\Gamma\right],
\end{align*}
这正是正文\eqnrefs{eq:fermionic_beam_coherent_incoherent_split}。

投影因子 $I-\Gamma$ 由费米反对易关系
$\hat c_j\hat c_k^\dagger=\delta_{jk}-\hat c_k^\dagger\hat c_j$
在 Wick 收缩中直接产生。若 $\Gamma=|u\rangle\langle u|$ 是单电子纯态，则
\[
\Tr[F^\dagger(I-\Gamma)F\Gamma]
=\langle u|F^\dagger F|u\rangle
-|\langle u|F|u\rangle|^2,
\]
正是单粒子驱动算符的方差。若 $\Gamma$ 是 $N_e$ 粒子 Slater 行列式的投影，$I-\Gamma$ 则把所有已占据末态从涨落跃迁中排除；这就是 Pauli 阻塞在一体驱动谱中的精确形式。

\end{proof}


为了把“相干电子列”与“随机独立到达”区分开，可以定义有限时间目标频率上的相干驱动幅
\begin{equation*}
\mathcal F_{ba}^{\rm coh}
\equiv
\int_{t_i}^{t_f}\dd t\,
\ee^{+\ii\omega_{ba}t}
\Tr[F(t)\Gamma]
\end{equation*}
以及涨落谱窗
\begin{equation*}
\mathcal S_{ba}^{\rm fl}
\equiv
\int\dd t\,\dd t'\,
\ee^{+\ii\omega_{ba}(t-t')}
\Tr\!\left[
F^\dagger(t')(I-\Gamma)F(t)\Gamma
\right].
\end{equation*}
于是正文\eqnrefs{eq:target_excitation_fermionic_beam_correlator} 分成
\[
P_{a\to b}^{(2)}
=\hbar^{-2}
\left(
|\mathcal F_{ba}^{\rm coh}|^2
+\mathcal S_{ba}^{\rm fl}
\right).
\]
这给出了电子束驱动的层级关系：平均电流只保留第一项；独立电子碰撞保留主要的 $O(N_e)$ 涨落项；相干整形束还可能出现 $O(N_e^2)$ 的相干项；简并或相关费米束则必须保留 Pauli/二体关联。Poisson 碰撞模型因此只适用于电子之间的相干和交换关联在相邻碰撞时间尺度上都不可分辨的参数区间。

\section{周期电子列与随机到达}

若电子逐个独立到达，并且每次碰撞之间靶系统经历同一个 CPTP 演化 $\mathscr E_0$，则一次周期的频闪映射是
\begin{equation}
\boxed{
\mathscr T_{\rm cyc}=\mathscr E_0\circ\mathscr E_e,
\qquad
\rho_A^{(n+1)}=\mathscr T_{\rm cyc}[\rho_A^{(n)}].}
\label{eq:feberi_stroboscopic_map}
\end{equation}
稳态满足
\begin{equation}
\boxed{
\rho_A^*=\mathscr T_{\rm cyc}[\rho_A^*].}
\label{eq:general_collision_fixed_point}
\end{equation}
求出固定点以后，还要检验从不同初态能否收敛到它。
\begin{proposition}[有限维频闪信道的收敛条件]
设 $\mathscr T$ 为有限维靶上的 CPTP 信道。它从每个初始密度算符收敛到同一个稳态 $\rho_*$，
当且仅当 $\ker(\mathscr T-I)=\mathbb C\rho_*$ 且所有其他本征值的模严格小于一。
若其他谱的最大模为 $r<1$，则对任意 $r<\eta<1$，存在有限常数 $C_\eta$ 使
$\|\mathscr T^n\rho-\rho_*\|_1\le C_\eta\eta^n$。
\end{proposition}
\begin{proof}
密度算符构成紧凸集，$\mathscr T$ 保持此集合，因而有固定态。
例如将 $n$ 次轨道的平均 $n^{-1}\sum_{j=0}^{n-1}\mathscr T^j\rho$ 取收敛子列，
其像与自身的差不超过 $2/n$，极限就是固定态。
CPTP 映射对厄米输入的迹范数不增；正负部分分解证明此结论，实部和虚部分解又给全算符空间中的一致有界性。
因此 $\mathscr T^n$ 不能在模一的本征值上有非平凡 Jordan 块，否则某个向量会随 $n$ 多项式增长。
在所述谱条件下，唯一不衰减分量是 $\rho_*\operatorname{Tr}\rho$；
其余 Jordan 块的 $n^k r^n$ 衰减对任意 $\eta>r$ 都有 $C_\eta\eta^n$ 上界，证明收敛及速率。
反过来，若每个初态收敛到同一态，由密度算符张成全部厄米算符、再作复线性延拓，
整个算符空间的极限必为 $X\mapsto\rho_*\operatorname{Tr}X$。
这排除额外的固定向量和其他模一本征值，得到必要性。
\end{proof}

常用的碰撞间演化可以由靶系统自身的 GKSL 生成元参数化：
\begin{equation}
\boxed{
\mathcal L_0\rho
=-\frac{\ii}{\hbar}[H_A,\rho]
+\sum_r\Gamma_r\mathcal D[L_r]\rho,
\qquad
\mathscr E_0=\ee^{T\mathcal L_0}.}
\label{eq:general_target_intercollision_lindblad}
\end{equation}
这里 $\Gamma_r\ge0$ 是衰变率，$L_r$ 是靶空间上的无量纲算符，
$\mathcal D[L]\rho=L\rho L^\dagger-\tfrac12\{L^\dagger L,\rho\}$。
$T$ 是相邻电子的锁定间隔；随机到达时，不能用平均间隔直接代替全部等待时间。若未追踪环境的记忆跨越多个碰撞，则不能复用同一个 $\mathscr T_{\rm cyc}$。此时应把长寿命环境模并入显式系统--记忆 Hilbert 空间，或使用\chapref{chap:prop-measure}\eqnrefs{eq:process_tensor_multitime_definition,eq:process_tensor_choi_link} 的多时刻过程张量。

周期到达需要逐次迭代同一个映射；随机到达则需要对碰撞次数和等待时间一起平均。
这一步可以对精确单次信道完成，并不要求每次作用很弱。
\begin{theorem}[独立随机碰撞的精确生成元]
设靶空间有限维，单次碰撞为固定 CPTP 信道 $\mathscr E_e$，
到达计数是速率 $R_e$ 的 Poisson 过程，碰撞可视为瞬时且与靶态和其他电子独立。
碰撞之间由固定 GKSL 生成元 $\mathcal L_0$ 演化。平均靶态满足
\begin{equation}
\dot\rho_A=\mathcal L_0\rho_A+R_e(\mathscr E_e-I)\rho_A.
\label{eq:collision-poisson-exact-generator}
\end{equation}
若 $K_a$ 是单次信道的 Kraus 算符，则碰撞部分已具有精确 GKSL 形式
\[
R_e(\mathscr E_e-I)=\sum_a\mathcal D[\sqrt{R_e}K_a].
\]
不要求 $\mathcal L_0$ 与 $\mathscr E_e$ 对易；非零碰撞时长、电子间关联及长寿命环境记忆须另行保留。
\end{theorem}
\begin{proof}
在 $\dd t$ 内，零次、一次、至少两次到达的概率分别为
$1-R_e\dd t+O(\dd t^2)$、$R_e\dd t+O(\dd t^2)$、$O(\dd t^2)$。
在零次路径上用 $I+\dd t\mathcal L_0$ 传播；一次碰撞发生在区间中哪个时刻，
只改变碰撞前后自由传播的 $O(\dd t)$ 修正，乘一次事件概率后成为 $O(\dd t^2)$。
于是
\[
\rho(t+\dd t)=(1-R_e\dd t)(I+\dd t\mathcal L_0)\rho(t)
+R_e\dd t\mathscr E_e\rho(t)+O(\dd t^2).
\]
相减除以 $\dd t$ 并取极限即得方程，整个推导不交换两个生成过程的次序。
又由 $\sum_aK_a^\dagger K_a=I$，
\[
\sum_a\mathcal D[\sqrt{R_e}K_a]\rho
=R_e\sum_aK_a\rho K_a^\dagger-\frac{R_e}{2}\{I,\rho\}
=R_e(\mathscr E_e-I)\rho.
\]
若 $\mathcal L_0=0$，对次数直接求和也给
\[
\rho(t)=\ee^{-R_et}\sum_{n=0}^\infty\frac{(R_et)^n}{n!}\mathscr E_e^n\rho(0)
=\ee^{R_et(\mathscr E_e-I)}\rho(0),
\]
其中每一项都是正权重的 CPTP 演化。
\end{proof}

\begin{exbox}[同一二能级靶的精确泵浦、稳态与退相干]
靶态为 $|g\rangle,|e\rangle$，能隙 $\hbar\omega_A$。
选择两个正交电子通道 $|h\rangle,|l\rangle$，其能量差同为 $\hbar\omega_A$，
每次入射制备为 $|h\rangle$。这是一个已截断并规定正交通道的模型。
令单次交换为
\[
U_\theta=\exp[-\ii\theta(|l,e\rangle\langle h,g|+|h,g\rangle\langle l,e|)],
\qquad 0\le\theta\le\pi/2.
\]
电子以速率 $R_e$ 独立随机到达，靶另有自发衰变率 $\gamma$。
求单次信道、读取出射电子后的条件靶态、平均稳态和相干衰减，并判断能否只用泵浦率描述碰撞。

\textbf{解。} 联合算符只旋转 $|h,g\rangle,|l,e\rangle$ 所张成的平面，
$U_\theta|h,g\rangle=c|h,g\rangle-\ii s|l,e\rangle$，
$U_\theta|h,e\rangle=|h,e\rangle$，其中 $c=\cos\theta$、$s=\sin\theta$。
在基底 $(g,e)$ 中，出射电子的两个 Kraus 算符为
\[
K_h=\begin{pmatrix}c&0\\0&1\end{pmatrix},\qquad
K_l=-\ii s|e\rangle\langle g|,\qquad K_h^\dagger K_h+K_l^\dagger K_l=I.
\]
若探测到低能电子，概率为 $p_l=s^2\rho_{gg}$，条件靶态为 $|e\rangle\langle e|$。
若探测到高能电子，概率为 $p_h=c^2\rho_{gg}+\rho_{ee}$，条件态为
$K_h\rho K_h^\dagger/p_h$。只有当相应概率非零时才定义条件态。
将两种未归一条件态相加，便恢复不读取电子时的信道；高能结果会重加权靶的两个能级，
故它也不是完全没有作用的事件。
令 $q=s^2$。单次通过使
$\rho_{ee}'=\rho_{ee}+q\rho_{gg}$，$\rho_{ge}'=c\rho_{ge}$。
定义 $\sigma_+=|e\rangle\langle g|$、$\sigma_-=\sigma_+^\dagger$、
$\sigma_z=|e\rangle\langle e|-|g\rangle\langle g|$。
由于 $K_h=(1+c)I/2+(1-c)\sigma_z/2$ 且两个系数为实数，恒等项的交叉贡献相消，故
\[
R_e(\mathscr E_e-I)
=R_eq\mathcal D[\sigma_+]+\frac{R_e(1-c)^2}{4}\mathcal D[\sigma_z].
\]
第一项泵浦，第二项还使相干额外衰减。
加入 $H_A=\hbar\omega_A|e\rangle\langle e|$ 和 $\gamma\mathcal D[\sigma_-]$ 后，
\begin{align*}
\dot n_e&=R_eq(1-n_e)-\gamma n_e,\\
\dot\rho_{ge}&=\left[\ii\omega_A-\frac\gamma2-R_e(1-c)\right]\rho_{ge}.
\end{align*}
若 $R_eq+\gamma>0$，人口的精确解和稳态为
\[
n_e(t)=\frac{R_eq}{R_eq+\gamma}
+\left[n_e(0)-\frac{R_eq}{R_eq+\gamma}\right]\ee^{-(R_eq+\gamma)t}.
\]
相干按 $\ee^{\ii\omega_At-[\gamma/2+R_e(1-c)]t}$ 演化。
例如 $\theta=\pi/2$、$R_e=\gamma$ 时，$n_e^*=1/2$，人口衰减率为 $2\gamma$，
相干衰减率为 $3\gamma/2$；只保留泵浦与自发衰变会误给 $\gamma$ 的相干衰减率。

在弱碰撞极限 $\theta\to0$，$q=\theta^2+O(\theta^4)$，
额外退相干的系数为 $R_e\theta^4/16+O(R_e\theta^6)$。
保持 $R_e\theta^2=\Gamma_p$ 固定，才得到仅含 $\Gamma_p\mathcal D[\sigma_+]$ 的极限泵浦模型。
固定间隔的电子列则应迭代 $\mathscr E_0\circ\mathscr E_e$；它的稳态与衰减不能由平均到达率直接代替。
在同一模型中，设电子严格每隔 $T$ 到达，碰撞后按 $\gamma$ 自发衰变。
记 $\lambda=\ee^{-\gamma T}$，在每次碰撞前读出的人口满足
\[
n_{k+1}=\lambda[q+(1-q)n_k],\qquad
n_{\rm pre}^*=\frac{\lambda q}{1-\lambda(1-q)}.
\]
上述分母须严格为正；若 $q=0$ 且 $\gamma=0$，人口保持初值，没有唯一人口稳态。
周期内的稳态时间平均为
$\bar n^*=q(1-\lambda)/[\gamma T(1-\lambda(1-q))]$，$\gamma T=0$ 时取连续极限。
例如 $q=1$、$\gamma T=1$，碰撞前人口是 $\ee^{-1}$，周期平均为 $1-\ee^{-1}\simeq0.6321$；
相同平均速率 $R_e=1/T$ 的随机束流稳态却是 $1/2$。
这个区别来自等待时间的统计，而不是改变了单次交换强度。
\end{exbox}

\begin{proposition}[二能级弱碰撞近似的有限时间误差]
在前一模型中，令 $\mathcal L$ 为精确平均生成元，
$\mathcal L_{\rm w}=\mathcal L_0+R_e\theta^2\mathcal D[\sigma_+]$。
对任意不参与碰撞的参考系统 $R$ 及归一化联合态 $\rho_{RA}$，
\[
\|[I_R\otimes(\ee^{t\mathcal L}-\ee^{t\mathcal L_{\rm w}})]\rho_{RA}\|_1
\le\min\!\left(2,\frac{19}{24}R_et\theta^4\right),\qquad t\ge0.
\]
因此传播时间也必须进入弱碰撞判据。
\end{proposition}
\begin{proof}
在 $0\le\theta\le\pi/2$，
$0\le\theta^2-\sin^2\theta\le\theta^4/3$，$0\le1-\cos\theta\le\theta^2/2$。
前一式可由 $\sin\theta\ge\theta-\theta^3/6\ge0$ 平方得到。
对迹类算符 $X$，次乘性给
$\|\mathcal D[L]X\|_1\le2\|L\|^2\|X\|_1$，参考扩张后同样成立。
所以
\[
\|(I_R\otimes(\mathcal L-\mathcal L_{\rm w}))X\|_1
\le\left(\frac{2R_e\theta^4}{3}+\frac{R_e\theta^4}{8}\right)\|X\|_1.
\]
两种传播都是 CPTP；对厄米 $X$，把它分成正负部分即可证明迹范数不增。
将上界代入两个传播子的 Duhamel 积分，得到 $19R_et\theta^4/24$。
两个归一化态的迹距离又不超过二，合并两界即得结论。
\end{proof}

\section{两电子的交换与可操作关联}

上面的电子束一体关联公式仍然把“电子之间是否能在第一量子化中区分”压缩进二阶量子化的一体密度矩阵。若两个自由电子的波包在时空上发生明显重叠，或者实验要直接层析电子--电子纠缠，就必须显式保留费米反对称性。单电子 Hilbert 空间记为 $\mathcal H_e$，则严格两电子空间不是
$\mathcal H_e\otimes\mathcal H_e$，而是其反对称子空间
\[
\wedge^2\mathcal H_e
=\mathsf A_-\bigl(\mathcal H_e\otimes\mathcal H_e\bigr),
\qquad
\mathsf A_-=\frac12(\mathbb I-\mathsf P_{12}),
\]
其中 $\mathsf P_{12}$ 交换两个电子的全部空间、自旋和其它内部自由度。

对两个已归一但不必正交的单电子波包 $|\phi\rangle,|\psi\rangle$，定义重叠
\[
s\equiv\langle\phi|\psi\rangle .
\]
只要 $|s|<1$，规范化反对称两电子态为
\begin{equation}
\boxed{
|\Phi_-\rangle
=
\frac{
|\phi\rangle_1|\psi\rangle_2
-
|\psi\rangle_1|\phi\rangle_2
}{
\sqrt{2(1-|s|^2)}
}.}
\label{eq:two_electron_antisym_wavepacket_state}
\end{equation}
重叠决定反对称态的归一化，但仅有小重叠还不能控制每个观测量的交换修正。记整个相互作用窗口内的最大重叠为
\[
\epsilon_{\rm ov}
\equiv
\sup_t|\langle\phi(t)|\psi(t)\rangle|
\ll1,
\]
还须对实际相互作用核和探测算符检查交叉矩阵元。共同的单粒子幺正传播保持内积，因此这个重叠本身也不能代替对局域交叉流的检查。

对任意交换对称的一体和两体算符
\[
\hat O
=
\hat o(1)+\hat o(2)+\hat W(1,2),
\qquad
[\hat O,\mathsf P_{12}]=0,
\]
反对称态矩阵元一般包含 直接与交换 两部分。尤其对环境积分掉以后得到的双电子二次核
\[
\hat W_K
=
\int\dd1\,\dd2\,
\hat j^\mu(1)\,
K_{\mu\nu}(1,2)\,
\hat j^\nu(2),
\]
这里要求实际构成的 $W_K$ 厄米且与交换算符对易。
核可来自推迟响应、噪声或低能约化，但原始推迟 Green 张量本身不必构成厄米能量算符，
须按所求观测量取相应核组合。在这些条件下，两电子矩阵元为
\begin{equation}
\boxed{
\langle\Phi_-|\hat W_K|\Phi_-\rangle
=
\frac{
W_{\rm dir}-\operatorname{Re}W_{\rm ex}^{(0)}
}{
1-|s|^2},
}
\label{eq:two_electron_direct_exchange_kernel}
\end{equation}
其中
\begin{align*}
W_{\rm dir}
&\equiv
\langle\phi\psi|\hat W_K|\phi\psi\rangle,\\
W_{\rm ex}^{(0)}
&\equiv
\langle\phi\psi|\hat W_K|\psi\phi\rangle.
\end{align*}
对正交波包 $s=0$，归一化分母消失，但 交换核本身未必自动为零；只有当算符核与两波包的时空支持使交叉矩阵元也很小，才真正回到“两个可区分电子各自沿一条轨迹”的近似。

\begin{proposition}[反对称双波包的直接与交换矩阵元]
设归一化单粒子波包 $\phi,\psi$ 的重叠 $s$ 满足 $|s|<1$，两体算符 $W$ 厄米 且与交换算符对易，各矩阵元存在。则规范化反对称态的期望为\eqnrefs{eq:two_electron_direct_exchange_kernel}。正交性 $s=0$ 不蕴含任意算符的交叉局域矩阵元为零。
\end{proposition}
\begin{proof}
先计算未归一反对称态
\[
|\widetilde\Phi_-\rangle
=
|\phi\psi\rangle-|\psi\phi\rangle
\]
的范数。利用
\[
\langle\phi\psi|\psi\phi\rangle
=
\langle\phi|\psi\rangle
\langle\psi|\phi\rangle
=
|s|^2,
\]
有
\begin{align*}
\langle\widetilde\Phi_-|\widetilde\Phi_-\rangle
&=
\langle\phi\psi|\phi\psi\rangle
+\langle\psi\phi|\psi\phi\rangle\\
&\quad-
\langle\phi\psi|\psi\phi\rangle
-\langle\psi\phi|\phi\psi\rangle\\
&=2(1-|s|^2),
\end{align*}
得到正文\eqnrefs{eq:two_electron_antisym_wavepacket_state}。

现在取任意满足
\[
[\hat W_K,\mathsf P_{12}]=0
\]
的厄米两体核。把规范化态代入矩阵元：
\begin{align*}
\langle\Phi_-|\hat W_K|\Phi_-\rangle
&=
\frac1{2(1-|s|^2)}
\Big[
\langle\phi\psi|\hat W_K|\phi\psi\rangle
+\langle\psi\phi|\hat W_K|\psi\phi\rangle\\
&\qquad
-\langle\phi\psi|\hat W_K|\psi\phi\rangle
-\langle\psi\phi|\hat W_K|\phi\psi\rangle
\Big].
\end{align*}
交换对称性给
\[
\langle\psi\phi|\hat W_K|\psi\phi\rangle
=
\langle\phi\psi|\hat W_K|\phi\psi\rangle
\equiv W_{\rm dir}.
\]
厄米性又给
\[
\langle\psi\phi|\hat W_K|\phi\psi\rangle
=
\langle\phi\psi|\hat W_K|\psi\phi\rangle^*.
\]
因此定义
\[
W_{\rm ex}^{(0)}
\equiv
\langle\phi\psi|\hat W_K|\psi\phi\rangle,
\]
便有
\[
\langle\Phi_-|\hat W_K|\Phi_-\rangle
=
\frac{
W_{\rm dir}-\operatorname{Re}W_{\rm ex}^{(0)}
}{
1-|s|^2},
\]
即正文\eqnrefs{eq:two_electron_direct_exchange_kernel}。
若只保留交换对称性而不要求厄米性，两条交叉矩阵元由交换对称性相等，
一般结果是 $(W_{\rm dir}-W_{\rm ex}^{(0)})/(1-|s|^2)$，可以为复数。
因而不能把任意因果响应核代入上面的实能量表达。

对按两个粒子槽分别作用的双点电流项
\[
\hat W_K^{(12)}
\equiv
\int\dd1\,\dd2\,
\hat j_1^\mu(1)K_{\mu\nu}(1,2)\hat j_2^\nu(2),
\]
定义交叉跃迁流
\[
j_{\phi\psi}^{\mu}(x)
\equiv
\langle\phi|\hat j^\mu(x)|\psi\rangle.
\]
张量积矩阵元逐槽因子化，于是交换矩阵元严格等于
\[
\langle\phi\psi|\hat W_K^{(12)}|\psi\phi\rangle
=
\int\dd1\,\dd2\,
j_{\phi\psi}^{\mu}(1)
K_{\mu\nu}(1,2)
j_{\psi\phi}^{\nu}(2).
\]
若完整两体核还含 $1\leftrightarrow2$ 的对称项，只需把交换后的同型积分相加。因此 $s=0$ 只说明两个波包在 Hilbert 内积意义下正交，并不单独保证所有局域交叉流都逐点为零。若两个波包在相互作用区的时空支持几乎不重叠，则 $j_{\phi\psi}$ 同时受抑，交换核才进一步趋零。这个额外条件正是把严格费米二电子问题约化为可标记电子轨迹模型所需的物理假设。
\end{proof}


这点对共享纳米光子环境尤其重要。两条可标记电子路径的近似需要在所保留的观测量上忽略交换矩阵元，而不只是让内积很小。若该条件失败，环境介导相位须与费米交换同时处理，实验层析对象应直接取
$\rho_{ee}\in L(\wedge^2\mathcal H_e)$。2026 年两电子量子行走实验已经需要在电子对的相干、经典关联、交换结构和纠缠之间作区分\cite{Tziperman2026TwoElectron}；这里给出的反对称连续波包一般表达式说明这种区分并不依赖能量格或路径量子比特表示。

\begin{exbox}[正交波包仍可有有限交换能]
在周期区间 $[0,L]$ 上取相同自旋的两个归一轨道
$\phi(x)=L^{-1/2}$、$\psi(x)=L^{-1/2}\ee^{\ii kx}$，$k=2\pi/L$。
设有界对称相互作用为 $W(x_1,x_2)=W_0\cos[k(x_1-x_2)]$。
计算直接、交换及反对称态的相互作用能，判断轨道正交是否足以舍去交换。

\textbf{解。} 一个周期积分给 $s=\langle\phi|\psi\rangle=0$。
两轨道的密度都为 $1/L$，所以
\[
W_{\rm dir}=\frac{W_0}{L^2}\int_0^L\int_0^L
\cos[k(x_1-x_2)]\dd x_1\dd x_2=0.
\]
交换矩阵元还含交叉相位，故
\[
W_{\rm ex}=\frac{W_0}{L^2}\int_0^L\int_0^L
\cos[k(x_1-x_2)]\ee^{\ii k(x_1-x_2)}\dd x_1\dd x_2=\frac{W_0}{2}.
\]
最后一步用 $\cos u\,\ee^{\ii u}=(1+\ee^{2\ii u})/2$，振荡项的周期积分为零。
因而 $\langle\Phi_-|W|\Phi_-\rangle=-W_0/2$。
这两个轨道内积为零，却遍布同一个空间区域，局域交叉密度并不消失。
是否可以标记两条电子路径，应由实际相互作用与探测核的交换矩阵元判断。
\end{exbox}

\subsection{Slater 分解与纯态关联熵}

对不可区分费米子，“把电子 1 对电子 2 做偏迹再算熵”还会遇到另一个概念问题：单个 Slater 行列式仅由反对称性就已经给出非零一体熵，不能把这部分自动解释为可利用纠缠。更稳妥的做法是先在有限实验工作窗中作两费米子 Slater 分解。任意归一纯态都可写成
\begin{equation}
\boxed{
|\Psi_{ee}\rangle
=\sum_{r=1}^{R_S}z_r\,
\hat f_{2r-1}^\dagger\hat f_{2r}^\dagger|0\rangle,
\qquad
\sum_r|z_r|^2=1,}
\label{eq:two_fermion_slater_decomposition}
\end{equation}
其中 $\{f_j\}$ 是一组正交单电子自然轨道，$R_S$ 是 Slater 秩。若 $R_S=1$，两电子态只是一个 Slater 行列式；若 $R_S>1$，才存在不能靠单粒子基变换消去的两费米子关联。

定义迹归一到 $1$ 的一体约化密度算符
\begin{equation}
\boxed{
\gamma^{(1)}_{ij}
\equiv
\frac12\langle\Psi_{ee}|\hat c_j^\dagger\hat c_i|\Psi_{ee}\rangle,
\qquad
\Tr\gamma^{(1)}=1.}
\label{eq:two_fermion_normalized_one_body_rdm}
\end{equation}
在 Slater 自然轨道基中，它的本征值成对出现：
\[
\lambda_{2r-1}=\lambda_{2r}=\frac{|z_r|^2}{2}.
\]
因此
\begin{equation}
\boxed{
S\!\left(\gamma^{(1)}\right)
=-\Tr\gamma^{(1)}\ln\gamma^{(1)}
=\ln2-\sum_r p_r\ln p_r,
\qquad p_r\equiv|z_r|^2,}
\label{eq:two_fermion_one_body_entropy_slater}
\end{equation}
其中不可避免的 $\ln2$ 是反对称单 Slater 行列式的交换基线。于是对纯两费米子态可定义
\begin{equation}
\boxed{
S_{\rm corr}
\equiv S(\gamma^{(1)})-\ln2
=-\sum_r p_r\ln p_r\ge0,}
\label{eq:two_fermion_correlation_entropy}
\end{equation}
它在且仅在 $R_S=1$ 时为零。这个量适合判断连续两电子纯态是否超出单 Slater 描述，但它不是所有混态和所有实验分区下唯一的纠缠度量。对混态或实验上把能量、路径、空间模式划成两个可操作子系统时，应在相应模式划分上使用完整 $\rho_{ee}$ 的负性、互信息或直接层析，而不能把纯态公式机械外推。

\begin{theorem}[纯两费米子态的 Slater 熵]
在有限维一粒子 Hilbert 空间中，任意规范化纯两费米子态具有\eqnrefs{eq:two_fermion_slater_decomposition}的正交 Slater 分解。令 $p_r=|z_r|^2$，迹一归一的一体密度本征值为每个 $p_r/2$ 各出现两次，故 $S(\gamma^{(1)})=\ln2-\sum_rp_r\ln p_r$，相关熵非负且仅在 Slater 秩一时为零。混态和不同可操作分区不能直接套用此纯态量。
\end{theorem}
\begin{proof}
先在有限实验工作窗中取正交单电子基 $\{|i\rangle\}_{i=1}^{d}$。任意纯两费米子态可写成
\[
|\Psi\rangle
=\frac12\sum_{ij}W_{ij}\hat c_i^\dagger\hat c_j^\dagger|0\rangle,
\qquad W^{\mathsf T}=-W.
\]
先构造酉合同所需的正交轨道对。若 $WW^\dagger u=s^2u$，$s>0$、$\|u\|=1$，令 $v=-W\bar u/s$。反对称性给 $u^\dagger v=0$，$\|v\|=1$，并且 $W\bar u=-sv$、$W\bar v=su$；其正交补在这个反线性作用下不变。逐对迭代并补上零空间基，便得到 $U^\dagger W\bar U=\bigoplus_r\left(\begin{smallmatrix}0&s_r\\-s_r&0\end{smallmatrix}\right)\oplus0$。这与正文 $f_\alpha^\dagger=\sum_iU_{i\alpha}c_i^\dagger$ 的轨道变换一致。

于是存在酉矩阵 $U$ 使
\[
U^\dagger W\bar U
=\bigoplus_{r=1}^{R_S}
\begin{pmatrix}0&z_r\\-z_r&0\end{pmatrix}\oplus0.
\]
定义 $\hat f_\alpha^\dagger=\sum_iU_{i\alpha}\hat c_i^\dagger$，便得到
\[
|\Psi\rangle
=\sum_rz_r\hat f_{2r-1}^\dagger\hat f_{2r}^\dagger|0\rangle.
\]
不同 $r$ 的行列式 正交，所以归一化要求 $\sum_r|z_r|^2=1$，得到正文\eqnrefs{eq:two_fermion_slater_decomposition}。

再计算归一化一体密度矩阵
\[
\gamma^{(1)}_{\alpha\beta}
=\frac12\langle\Psi|\hat f_\beta^\dagger\hat f_\alpha|\Psi\rangle.
\]
固定 $r$ 的两个自然轨道各占据一次，因此
\[
\langle\hat n_{2r-1}\rangle
=\langle\hat n_{2r}\rangle
=|z_r|^2,
\]
且自然轨道基中的非对角元为零。于是
\[
\gamma^{(1)}
=\bigoplus_r\frac{|z_r|^2}{2}
\begin{pmatrix}1&0\\0&1\end{pmatrix}.
\]
令 $p_r=|z_r|^2$，则
\begin{align*}
S(\gamma^{(1)})
&=-\sum_r2\frac{p_r}{2}\ln\frac{p_r}{2}\\
&=\ln2-\sum_rp_r\ln p_r.
\end{align*}
单一 Slater 行列式对应 $p_1=1$，仍有 $S=\ln2$；这只是反对称交换基线。故
\[
S_{\rm corr}=S-\ln2=-\sum_rp_r\ln p_r\ge0,
\]
且只有 $R_S=1$ 时为零，得到正文\eqnrefs{eq:two_fermion_correlation_entropy}。
连续波包的反对称 Hilbert--Schmidt 核可产生可数 Slater 分解，
有限秩截断可使一体密度按迹范数收敛；但无限维熵还须有可求和的熵尾或相应能量约束，
不能由核的范数收敛直接推出熵收敛。

\end{proof}


\begin{exbox}[相同一体密度不决定两电子的量子关联]
取四个正交模式，记 $|ij\rangle=f_i^\dagger f_j^\dagger|0\rangle$，$i<j$。
比较 $|12\rangle$、$|\Psi\rangle=(|12\rangle+|34\rangle)/\sqrt2$ 及
$\rho_{\rm mix}=(|12\rangle\langle12|+|34\rangle\langle34|)/2$ 的一体熵。
再把模式分为 $A=\{1,3\}$、$B=\{2,4\}$，判断后两态在该可操作分区的关联。

\textbf{解。} 单一行列式的一体密度为
$\operatorname{diag}(1/2,1/2,0,0)$，所以 $S=\ln2$，纯态相关熵为零。
对 $|\Psi\rangle$，两个 Slater 权重均为 $1/2$，故
\[
\gamma^{(1)}=I_4/4,\qquad S=\ln4,\qquad S_{\rm corr}=\ln2.
\]
混合态 $\rho_{\rm mix}$ 给出完全相同的一体密度和一体熵，
但纯态相关熵公式的假设已不满足。

在题设分区中，每一项都在 $A,B$ 各有一个电子。
固定按 $A$ 后 $B$ 的模式顺序，$|\Psi\rangle$ 对应
$(|1\rangle_A|2\rangle_B+|3\rangle_A|4\rangle_B)/\sqrt2$，
是 Schmidt 权重均为 $1/2$ 的纠缠态；两侧约化熵均为 $\ln2$。
而 $\rho_{\rm mix}$ 是这两个乘积态的等权混合，属于可分态。
它的一体熵虽然也为 $\ln4$，却只包含交换基线及混合的不确定性。
因此测得同一自然轨道占据数，仍须读取两电子相干，才能区分这两个物理状态。
\end{exbox}

若两电子再与未探测环境耦合，物理末态为
\[
\rho_{ee}'
=\Tr_E\!\left[
U_{eeE}(\rho_{ee}\otimes\rho_E)U_{eeE}^\dagger
\right].
\]
只要总 Hamiltonian 与电子交换算符对易，演化始终保持在 $\wedge^2\mathcal H_e$ 中；环境会改变 Slater 权重和相干，却不会把费米反对称性“退相干掉”。因此实验上看到的两电子经典相关、交换相关和真正量子关联必须在同一个反对称密度算符中区分。

若照明场已经制备为已知的相干近场，电子还可以由受控的外场传播描述。\chapref{chap:external-pinem}将从这一具体制备出发，计算近场如何形成能量边带；其量子交换和反冲条件沿用前面已经建立的连续态与通道结构。
